\documentclass[a4paper,11pt]{article}

% set up sensible margins (same as for cssethesis)
\usepackage[paper=a4paper,left=30mm,width=150mm,top=25mm,bottom=25mm]{geometry}
\usepackage{setspace}               % This is used in the title page
\usepackage{graphicx}               % This is used to load the crest in the title page
\usepackage{poltakmacros}           % Personal macros included in file 'poltakmacros.sty'
\usepackage[font={small}]{caption}
\usepackage[hidelinks]{hyperref}
\usepackage{url}

\author{Jonathan Poltak Samosir}
\title{Honours Research Proposal}

\begin{document}

% Set up a title page
\thispagestyle{empty} % no page number on very first page
% Use roman numerals for page numbers initially
\renewcommand{\thepage}{\roman{page}}

\begin{spacing}{1.5}
\begin{center}
{\Large \bfseries
Clayton School of Information Technology\\
Monash University}

\vspace*{30mm}

\includegraphics[width=5cm]{img/MonashCrest.pdf}

\vspace*{15mm}

{\large \bfseries
Honours Research Proposal --- Semester 2, 2014
}

\vspace*{10mm}

{\LARGE \bfseries
A pipeline for the preprocessing and storage of heterogeneous big data [WORKING TITLE]
}

\vspace*{20mm}

{\large \bfseries
Jonathan Poltak Samosir [2271 3603]

\vspace*{20mm}

Supervisors: \parbox[t]{50mm}{\mbox{Dr Maria Indrawan-Santiago}\\Dr Pari Delir Haghighi}
}

\end{center}
\end{spacing}

\newpage

\tableofcontents

\newpage
% Now reset page number counter,and switch to arabic numerals for remaining
% page numbers 
\setcounter{page}{1}
\renewcommand{\thepage}{\arabic{page}}

\section{Introduction}
In this chapter I will demonstrate some of the extended citation
capabilities provided by the {\sf natbib} package \cite{Dal1999}. It
replaces the standard \LaTeX\ \verb+\cite+ command with two basic forms of
citation command: \verb+\citep+ and \verb+\citet+, as well as providing
several other very useful ones.

The \verb+\citep+ command is best used when placing a citation at the end
of sentence or phrase (as above).  In the {\sf natbib} documentation, this
is referred to as a \emph{parenthetical citation}.%
\footnote{For ease of conversion from exisiting \LaTeX\ documents, you
might find it useful to place
\texttt{$\backslash$renewcommand\{$\backslash$cite\}\{$\backslash$citep\}}
in the preamble of the document, since any existing standard
\texttt{$\backslash$cite} commands should almost certainly be treated as
\texttt{$\backslash$citep}.}

When you want to refer to the authors of a particular work, typically at
the start of a sentence, a parenthetical citation is not appropriate. This
is particularly so if you are using a numerical or symbolic citation style.
You should \emph{not} start a sentence with
\begin{quote}
[2] says that this is most certainly \ldots
\end{quote}
In such situations you really need to give the authors' names. The
\verb+\citet+ command produces \emph{textual citations}, which allows you
to produce things like:
\begin{quote}
\citet{Ade1983} describes a means by which textures may be
characterized \ldots another approach is given in \citet{DeV1998}.

\citet{AGR1996} note that humans have little or no difficulty in
perceiving shape, yet find it extremely difficult to \emph{describe} what
they perceive.
\end{quote}

Note that an abbreviated version of the authors' names has been produced in
the third example above.  It is often desirable to have the full list of
authors' names given when a work  is first cited, and an abbreviated list
thereafter. This can be achieved by passing the \texttt{longnamesfirst}
option to {\sf natbib} when the package is used. This will produce
an initial citation like:
\begin{quote}
\citet*{AGR1996} note that humans have little or no difficulty in
perceiving shape, yet find it extremely difficult to \emph{describe} what
they perceive.
\end{quote}

Both the \verb+\citep+ and \verb+\citet+ can take two optional arguments.
If just one is provided,  its text will appear as a ``post-note'' after the
citation details. If two arguments are provided, the first defines a
pre-note, and the second a post-note. Here is an example:
\begin{quote}
\verb+\citep[Ch.~3]{AaK1989}+ \ldots{} \citep[Ch.~3]{AaK1989}\\
\verb+\citep[see][Ch.~3]{AaK1989}+ \ldots{} \citep[see][Ch.~3]{AaK1989}\\
\end{quote}

These examples only scratch the surface of what the {\sf natbib} package
can do. To discover the full power of the package, see the documentation at
CTAN \citep{Dal1999}. You probably already have it on your system. Try
\verb+locate natbib.dvi+ at the command line.

\section{Research Context}
Duis autem vel eum iriure dolor in hendrerit in vulputate velit esse molestie
consequat, vel illum dolore eu feugiat nulla facilisis at vero eros et accumsan
et iusto odio dignissim qui blandit praesent luptatum zzril delenit augue duis
dolore te feugait nulla facilisi. Lorem ipsum dolor sit amet, consectetuer
adipiscing elit, sed diam nonummy nibh euismod tincidunt ut laoreet dolore
magna aliquam erat volutpat.

\section{Research Plan and Methods}
Lorem ipsum dolor sit amet, consetetur sadipscing elitr,  sed diam nonumy
\subsection{Methodologies}
eirmod tempor invidunt ut labore et dolore magna aliquyam erat, sed diam
voluptua. At vero eos et accusam et justo duo dolores et ea rebum. Stet clita
kasd gubergren, no sea takimata sanctus est Lorem ipsum dolor sit amet. Lorem
ipsum dolor sit amet, consetetur sadipscing elitr,  sed diam nonumy eirmod
tempor invidunt ut labore et dolore magna aliquyam erat, sed diam voluptua. At
vero eos et accusam et justo duo dolores et ea rebum. Stet clita kasd
gubergren, no sea takimata sanctus est Lorem ipsum dolor sit amet. Lorem ipsum
dolor sit amet, consetetur sadipscing elitr,  sed diam nonumy eirmod tempor
invidunt ut labore et dolore magna aliquyam erat, sed diam voluptua. At vero
eos et accusam et justo duo dolores et ea rebum. Stet clita kasd gubergren, no
sea takimata sanctus est Lorem ipsum dolor sit amet.

\section{Deliverables}
Duis autem vel eum iriure dolor in hendrerit in vulputate velit esse molestie
consequat, vel illum dolore eu feugiat nulla facilisis at vero eros et accumsan
et iusto odio dignissim qui blandit praesent luptatum zzril delenit augue duis
dolore te feugait nulla facilisi. Lorem ipsum dolor sit amet, consectetuer
adipiscing elit, sed diam nonummy nibh euismod tincidunt ut laoreet dolore
magna aliquam erat volutpat.


\bibliographystyle{dcu}
\bibliography{examplebib}

\end{document}
